\documentclass{article}

\usepackage[margin=1in]{geometry}
\linespread{1.2}
\usepackage{amsmath,amsthm,amssymb,amsfonts,commath}


\newcommand{\mto}{\stackrel{\mu}{\rightarrow}}
\newcommand{\Mcal}{\mathcal{M}}


 
\title{Real Analysis I --- Homework 5}

% Replace with your name and ID below
\author{FirstName LastName \\ ID: 123-45-6789}

\date{}


\begin{document}

\maketitle

Unless otherwise specified, all functions are defined on $E \subset \mathbb{R}^{n}$ where $E \in \mathcal{M}$.
\bigskip

For notation simplicity, a set such as $\{x \in E: f(x) > t\}$ is written as $\{f>t\}$ for short.

\begin{enumerate}

\item Suppose $f \in L(E)$. Let $E_{k}$ be non-decreasing and $\lim_{k} E_{k} = E$. Show that
\[
\lim_{k \to \infty} \int_{E_{k}} f = \int_{E} f .
\]

%uncomment below and write your proof after it.
%\begin{proof}
%\end{proof}


%\newpage % uncomment this line to start on a new page

\item Suppose $f \in L(E)$ where $\mu(E) < \infty$. Let $E_{k,j} = \{\frac{j}{k} \le f < \frac{j+1}{k}\}$. Show that
\[
\lim_{k \to \infty} \sum_{j=-\infty}^{\infty} \frac{j}{k}\cdot \mu(E_{k,j}) = \int_{E} f.
\]

%uncomment below and write your proof after it.
%\begin{proof}
%\end{proof}


%\newpage % uncomment this line to start on a new page

\item Suppose $f \in L(\mathbb{R})$ and $f(0) = 0$, and $f$ is differentiable at $x=0$. Show that $\frac{f(x)}{x} \in L(\mathbb{R})$.

%uncomment below and write your proof after it.
%\begin{proof}
%\end{proof}



%\newpage % uncomment this line to start on a new page

\item Let $f \in C([a,b])$ (i.e., $f: [a,b] \to \mathbb{R}$ is continuous). Show that
\[
\lim_{k\to\infty} \int_{a}^{b} f(x) \sin(kx) \dif x = 0. 
\]
Hint: For any $\epsilon>0$, consider a proper partition $\Delta: a=x_{0}< \dots < x_{n}=b$ and $K \in \mathbb{N}$ such that the absolute value of the right hand side of  
\[
\int_{a}^{b} f(x) \sin(kx) \dif x = \sum_{i=1}^{n} \int_{x_{i-1}}^{x_{i}}(f(x) - f(x_{i})) \sin(kx) \dif x + \sum_{i=1}^{n} \int_{x_{i-1}}^{x_{i}} f(x_{i}) \sin(kx)\dif x
\]
can be bounded by $\epsilon$ for all $k \ge K$.

%uncomment below and write your proof after it.
%\begin{proof}
%\end{proof}



%\newpage % uncomment this line to start on a new page

\item Use the fact in the previous problem to show that the claim is true if $f \in L([a,b])$.

%uncomment below and write your proof after it.
%\begin{proof}
%\end{proof}



%\newpage % uncomment this line to start on a new page

\item Suppose $\mu(E) < \infty$ and $f_{k}: E \to \mathbb{R}$ are measurable. Show that $f_{k} \mto 0$ on $E$ iff
\[
\lim_{k\to \infty }\int_{E} \frac{|f_{k}|}{1+|f_{k}|} = 0.
\]

%uncomment below and write your proof after it.
%\begin{proof}
%\end{proof}


%\newpage % uncomment this line to start on a new page


\item Suppose $f: E \to \mathbb{R}$, and for any $\epsilon > 0$, there exist $g,h \in L(E)$ such that $g \le f \le h$ and
\[
\int_{E} |h - g |\le \epsilon.
\]
Show that $f$ is measurable on $E$ and $f \in L(E)$.

%uncomment below and write your proof after it.
%\begin{proof}
%\end{proof}




%\newpage % uncomment this line to start on a new page


\item Suppose $f_{k},f: E \to \mathbb{R}$ are measurable, $f_{k} \to f$ a.e., and there exists $M>0$ such that
\[
\int_{E} |f_{k}| \le M
\]
for all $k$. Show that $f \in L(E)$.

%uncomment below and write your proof after it.
%\begin{proof}
%\end{proof}



%\newpage % uncomment this line to start on a new page


\item 
Let $F \subset [0,1]$ be closed and $\mu(F) = 0$. Show that $\chi_{F}$ is Riemann integrable on $[0,1]$.

%uncomment below and write your proof after it.
%\begin{proof}
%\end{proof}




%\newpage % uncomment this line to start on a new page


\item 
Suppose $f : [0,1] \to [a,b]$ is Riemann integrable, $g \in C([a,b])$. Show that $g\circ f$ is Riemann integrable on $[0,1]$.

%uncomment below and write your proof after it.
%\begin{proof}
%\end{proof}





\end{enumerate}


\end{document}